% Part: sets-functions-relations
% Chapter: arithmetization
% Section: checking-details
%
\documentclass[../../../include/open-logic-section]{subfiles}

\begin{document}

\olfileid{sfr}{arith}{check}
\olsection{ترتیب دار حلقے تے میدان}

ایس پورے باب وچ اسی دعویٰ کیتا کہ کجھ تعریفاں ”جیویں اوہناں نوں کرنا
چاہیدا اے“ اوویں ورتاؤ کردیاں نیں۔ ایس فنی ضمیمے وچ اسی کھول کے دساں
گے کہ ساڈا مطلب کی اے، تے (ایہدا خاکہ دیاں گے کہ) تعریفاں ”درست“
ورتاؤ کردیاں نیں۔

\olref[int]{sec} وچ اسی $\Int$ اُتے جمع تے ضرب تعریف کیتے سن۔ اسی
وکھاؤنا چاہنے آں کہ جیویں اوہ تعریف کیتے گئے نیں، اوہ $\Int$ نوں اوہ
ساخت دیندے نیں جیہڑی اسی ایہدے کول ”چاہواں گے“۔ خاص طور تے متعلقہ ساخت
تبادلی حلقے دی اے۔

\begin{defn}
	اک \emph{تبادلی حلقہ} اک سیٹ $S$ اے، جیہدے نال خاص رکن $0$ تے $1$
	تے عمل $+$ تے $\times$ دِتے گئے نیں، تے جیہڑا ایہ اَٹھ فارمولے پورے کردا اے:
	\begin{align*}
		\emph{شرکتیت}&&a + (b+ c) & = (a + b) + c \\
		&& (a \times b) \times c & = a \times (b\times c)\\
		\emph{تبادلیت}&&a + b &= b+ a  \\
		&&  a \times b&= b\times a\\
		\emph{شناختی رکن}&&a + 0 &= a \\
		&& a \times 1 &= a\\
		\emph{جمعی اُلٹ}&&(\exists b\in S)0&=a + b\\
		\emph{تقسیمیت}&&a \times (b+ c ) &= (a \times b) + (a \times c)
	\end{align*}
	لُکّویں طور تے ایہ سارے $S$ تک محدود کلی مقدارا نال بندھے ہوئے نیں۔
	تے یاد رکھو کہ ایتھے $0$ تے~$1$ ناں والے رکن لازمی نہیں کہ ایسے ناں
	والے قدرتی عدد ہون۔
\end{defn}

سو ایہ جانچن لئی کہ صحیح عدد اک تبادلی حلقہ بناؤندے نیں، سانوں صرف ایہ
جانچنا اے کہ اسی ایہ اَٹھ شرطاں پوری کردے آں۔ کوئی شرط ثابت کرنا
{اوکھا} نہیں، پر کم ذرا لمبا اے۔ مثال لئی جمع دے معاملے وچ
\emph{شرکتیت} ایویں ثابت ہوندی اے:

\begin{proof}
$i, j, k \in \Int$ مقرر کرو۔ سو کجھ $a_1, b_1, a_2, b_2, a_3, b_3 \in
\Nat$ نیں جیہناں لئی $i = \equivrep{a_1, b_1}{}$ تے $j =
\equivrep{a_2,b_2}{}$ تے $k = \equivrep{a_3, b_3}{}$۔ (پڑھن دی
سوکھ لئی اسی ”$\equivrep{x, y}{}$“ لکھدے آں، نہ کہ
”$\equivrep{x, y}{\Intequiv}$“؛ ایس پورے حصے وچ ایویں کراں گے۔) ہُن:
\begin{align*}
	i + (j + k) &= \equivrep{a_1, b_1}{}+(\equivrep{a_2, b_2}{} + \equivrep{a_3, b_3}{}) \\
	&= \equivrep{a_1,  b_1}{} + \equivrep{a_2+a_3, b_2+b_3}{}\\
	&= \equivrep{a_1 + (a_2 + a_3), b_1 + (b_2 + b_3)}{}\\
	&= \equivrep{(a_1 + a_2) + a_3, (b_1 + b_2) + b_3}{}\\
	&= \equivrep{a_1 + a_2, b_1 + b_2}{} + \equivrep{a_3, b_3}{}\\
	&= (\equivrep{a_1, b_1}{} + \equivrep{a_2, b_2}{}) + \equivrep{a_3, b_3}{}\\
	&= (i+j) + k
\end{align*}
ایتھے اسی $\Nat$ اُتے جمع دے ورتاؤ نوں کھُلھ کے ورتیا اے۔
\end{proof}

ایسے طرح \emph{جمعی اُلٹ} دا ثبوت ایہ اے:

\begin{proof}
$i \in \Int$ مقرر کرو، سو $i = \equivrep{a,b}{}$، کجھ $a,b \in
\Nat$ لئی۔ $j = \equivrep{b,a}{} \in \Int$ مقرر کرو۔
قدرتی عدداں دے ورتاؤ نوں ورت کے $(a+b) + 0 = 0 + (a+b)$، ایس لئی
تعریف دے مطابق $\tuple{a+b, b+a} \sim_\Int \tuple{0,0}$، تے نتیجتاً
$\equivrep{a+b, b+a}{} = \equivrep{0, 0}{} = 0_\Int$۔ سو ہُن $i + j =
\equivrep{a,b}{}+\equivrep{b,a}{}=\equivrep{a+b, b+a}{}= \equivrep{0,
0}{} = 0_\Int$۔
\end{proof}

تے \emph{تقسیمیت} دا ثبوت ایہ اے:

\begin{proof}
اوپر وانگ $i = \equivrep{a_1, b_1}{}$ تے $j =
\equivrep{a_2,b_2}{}$ تے $k = \equivrep{a_3, b_3}{}$ مقرر کرو۔ ہُن:
\begin{align*}
	i \times (j + k)
	&= \equivrep{a_1, b_1}{} \times (\equivrep{a_2,b_2}{} + \equivrep{a_3, b_3}{})\\
	&= \equivrep{a_1, b_1}{} \times \equivrep{a_2 + a_3,b_2+b_3}{}\\
	&= \equivrep{a_1  (a_2 + a_3) + b_1  (b_2+b_3), a_1  (b_2 + b_3) + b_1 (a_2 + a_3)}{}\\
	&= \equivrep{a_1 a_2 + a_1a_3 + b_1 b_2+b_1b_3, a_1 b_2 + a_1b_3 + a_2b_1 + a_3b_1}{}\\
%		&= \equivrep{a_1 a_2 + b_1b_2 + a_1a_3 + b_1 b_2+b_1b_3, a_1 b_2 + a_1b_3 + a_2b_1 + a_3b_1}{}\\
	&= \equivrep{a_1a_2 + b_1b_2, a_1b_2 + a_2b_1}{} + \equivrep{a_1a_3 + b_1b_3, a_1b_3 + a_3b_1}{}\\
	&= (\equivrep{a_1, b_1}{} \times \equivrep{a_2,b_2}{}) + (\equivrep{a_1, b_1}{} \times  \equivrep{a_3, b_3}{})\\
	&= (i \times j) + (i \times k)
\end{align*}
\end{proof}

باقی پنج شرطاں دا ثبوت اسی مشق لئی چھڈدے آں۔ اوہ کر لین توں بعد اسی
وکھا دِتا اے کہ $\Int$ اک تبادلی حلقہ بناؤندا اے، یعنی جمع تے ضرب
(جیویں اوہ تعریف کیتے گئے نیں) اوویں ورتاؤ کردے نیں جیویں اوہناں نوں
کرنا چاہیدا اے۔

\begin{prob}
ثابت کرو کہ $\Int$ اک تبادلی حلقہ اے۔
\end{prob}

پر ساڈا کم ہُنّے پورا نہیں ہویا۔ $\Int$ اُتے جمع تے ضرب دے نال نال
اسی اک ترتیبی تعلق $\leq$ وی تعریف کیتا سی، تے سانوں جانچنا اے کہ ایہ
وی اوویں ورتاؤ کردا اے جیویں ایہنوں کرنا چاہیدا اے۔ ہور تفصیل نال سانوں
وکھاؤنا اے کہ $\Int$ اک \emph{ترتیب دار} حلقہ بناؤندا اے۔\footnote{
	\olref[sfr][rel][ord]{def:linearorder} توں یاد کرو کہ کُل ترتیب اوہ
	تعلق اے جیہڑا انعکاسی، متعدی، یک طرفہ تے جڑیا ہووے۔ ترتیبی تعلقاں دے
	سیاق وچ جڑاؤ نوں کدی کدی \emph{سہ شقیت} آکھیا جاندا اے، کیوں جو ہر
	$a$ تے $b$ لئی ساڈے کول $a \leq b \lor a
	= b \lor a \geq b$ اے۔}

\begin{defn}
اک \emph{ترتیب دار حلقہ} اوہ تبادلی حلقہ اے جیہدے نال اک کُل ترتیبی
تعلق $\leq$ وی دِتا گیا اے، تے جیہڑا ایہ شرطاں پوری کردا اے:
\begin{align*}
	a \leq b &\lif a + c \leq b + c\\
	(a \leq b \land 0 \leq c) &\lif a \times c \leq b \times c
\end{align*}
\end{defn}

\begin{prob}
ثابت کرو کہ $\Int$ اک ترتیب دار حلقہ اے۔
\end{prob}

پہلے وانگ، ایہ وکھاؤنا لمبا پر معمول دا کم اے کہ بناوٹی $\Int$ اک
ترتیب دار حلقہ اے۔ اسی ایہ کم تہاڈے لئی چھڈدے آں۔

ایس نال صحیح عدداں دا کم پورا ہو گیا۔ پر ہُن سانوں ناطق عدداں بارے بہت
رلدیاں گلاں وکھاؤنیاں نیں۔ خاص طور تے $+$، $\times$ تے $\leq$ دیاں
ساڈیاں دِتیاں تعریفاں تحت ہُن سانوں وکھاؤنا اے کہ ناطق عدد اک ترتیب دار
\emph{میدان} بناؤندے نیں:
\begin{defn}\ollabel{orderedfield}
اک \emph{ترتیب دار میدان} اوہ ترتیب دار حلقہ اے جیہڑا ایہ شرط وی پوری کردا اے:
\begin{align*}
	\emph{ضربی اُلٹ}& & (\forall a \in S \setminus \{0\})(\exists b \in S) a\times b& = 1
\end{align*}
\end{defn}

جدوں تسیں وکھا دِتا کہ $\Int$ اک ترتیب دار حلقہ بناؤندا اے، تاں ایہ
وکھاؤنا سوکھا پر لمبا اے کہ $\Rat$ اک ترتیب دار میدان بناؤندا اے۔

\begin{prob}
ثابت کرو کہ $\Rat$ اک ترتیب دار میدان اے۔
\end{prob}

صحیح تے ناطق عدداں نال نمٹ کے صرف حقیقی عدد رہ جاندے نیں۔ خاص طور تے
سانوں وکھاؤنا اے کہ $\Real$ اک \emph{مکمل} ترتیب دار میدان بناؤندا اے،
یعنی تکمیلیت دی خاصیت والا ترتیب دار میدان۔ ہُن
\olref[cuts]{realcompleteness} نے ثابت کیتا کہ $\Real$ کول تکمیلیت دی
خاصیت اے۔ پر ہُنّے ایہ جانچن دا (اُکتاؤنا) کم کرنا باقی اے کہ $\Real$
اک ترتیب دار میدان اے۔
% PNBARI-008 ماخذ دی درستی: انگریزی جملے وچ ”the (tedious) of checking“ دے وچکار کم والا ناں رہ گیا سی؛ ایتھے صرف جانچن دا کم آکھیا اے۔

اوس لمبی مشق ول \emph{دوڑن} توں پہلاں سانوں کجھ ہور ”فوری“ گلاں
جانچنیاں پینیاں نیں۔ مثال لئی سانوں ضمانت چاہیدی اے کہ $\alpha +
\beta$، جیویں ایہ تعریف کیتا گیا اے، ہر کٹ $\alpha$ تے $\beta$ لئی
واقعی اک \emph{کٹ} اے۔ ایس حقیقت دا ثبوت ایہ اے:

\begin{proof}
کیوں جو $\alpha$ تے $\beta$ دونیں کٹ نیں، $\alpha + \beta = \Setabs{p
+ q}{p \in \alpha \land q \in \beta}$، $\Rat$ دا اک غیر خالی حقیقی
ذیلی سیٹ اے۔ ہُن فرض کرو کہ $x < p + q$، کجھ $p \in \alpha$ تے
$q \in \beta$ لئی۔ فیر $x - p < q$، سو $x - p \in \beta$، تے
$x = p + (x - p) \in \alpha + \beta$۔ سو $\alpha + \beta$، $\Rat$
دا اک مُڈھلا ٹکڑا اے۔ آخر وچ ہر $p + q \in \alpha + \beta$ لئی، کیوں
جو $\alpha$ تے $\beta$ دونیں کٹ نیں، کجھ $p_1 \in \alpha$ تے
$q_1 \in \beta$ نیں جیہناں لئی $p < p_1$ تے $q < q_1$؛ سو
$p + q < p_1 + q_1 \in \alpha + \beta$؛ سو $\alpha + \beta$ دا کوئی
سب توں وڈا رکن نہیں۔
\end{proof}

ایہو جیہیاں کوششاں نال تسیں جانچ سکدے او کہ $\alpha - \beta$ تے
$\alpha \times \beta$ تے $\alpha \div \beta$ کٹ نیں (آخری معاملے وچ
$\beta$ دے صفر-کٹ ہون والے معاملے نوں چھڈ کے)۔ پر اک واری فیر اسی ایہ
کم تہاڈے لئی چھڈدے آں۔
% PNBARI-012 ماخذ دی وضاحت: cuts.tex وچ تفریق تے تقسیم دیاں دکھائیاں تعریفاں غیر فعال تبصری سطراں نیں؛ ایتھے جانچ دی درخواست محفوظ اے، پر اوہ تعریفاں فعال نہیں کیتیاں گئیاں۔

\begin{prob}
ثابت کرو کہ $\Real$ اک ترتیب دار میدان اے۔
\end{prob}

پر ایتھے اک چھوٹا کھلا سرا سمیٹنا اے۔ \olref[cuts]{sec} وچ اسی
سجھایا سی کہ $\sqrt{2} =
\Setabs{p \in \Rat}{p < 0 \text{ یا }p^2 < 2}$ لیا جا سکدا اے۔ پر سانوں
ایہ وکھاؤنا پینا اے کہ ایہ سیٹ اک \emph{کٹ} اے۔ ایس حقیقت دا ثبوت ایہ اے:

\begin{proof}
صاف اے کہ ایہ ناطق عدداں دا اک غیر خالی حقیقی مُڈھلا ٹکڑا اے؛ سو اینا
وکھاؤنا کافی اے کہ ایہدا کوئی سب توں وڈا رکن نہیں۔ خاص طور تے اینا
وکھاؤنا کافی اے کہ جتھے $p$ اک مثبت ناطق عدد اے جیہدے لئی $p^2 < 2$
تے $q = \frac{2p+2}{p+2}$، اوتھے $p < q$ تے $q^2 < 2$ دونیں۔
$p < q$ ویکھن لئی صرف ایہ نوٹ کرو:
\begin{align*}
	p^2 &< 2\\
	p^2 + 2p &< 2 + 2p\\
	p(p + 2) &< 2 + 2p\\
	p &< \tfrac{2+2p}{p+2} = q
\end{align*}
$q^2 < 2$ ویکھن لئی صرف ایہ نوٹ کرو:
\begin{align*}
	p^2 &< 2\\
	2p^2 + 4p + 2 &< p^2 + 4p+ 4\\
	4p^2 + 8p + 4 &< 2(p^2 + 4p + 4)\\
	(2p+2)^2 & <2(p+2)^2\\
	\tfrac{(2p+2)^2}{(p+2)^2} &< 2\\
	q^2 &< 2
\end{align*}
\end{proof}
\end{document}
